% !TeX TS-program = xelatex
\PassOptionsToClass{fontset=fandol}{ctexart}
\documentclass{resume}
\ResumeName{陈沫樵}
\usepackage{enumitem}

\begin{document}

\ResumeContacts{
  (+86) 186\-1195\-1800,%
  \ResumeUrl{mailto:mochiaochen@gmail.com}{mochiaochen@gmail.com},%
  \ResumeUrl{https://github.com/MochiaoChen}{github.com/MochiaoChen}
}

\ResumeTitle

\section{学校}
\ResumeItem
[对外经济贸易大学 | 本科]
{\textbf{对外经济贸易大学}}
[{\zihao{5} 金融学（鸿儒金融人才培养实验班），中国金融学院 | 经济学学士}]
[2024.09—2028.06]

\section{实习}

\ResumeItem{\textbf{AI 应用与评测实习生}}
[UniPat AI]
[2026.08—2026.09]
\begin{itemize}
  \item \textbf{AI 质检产品设计}：针对专家审核中条件与示例错配、原文查找困难、反馈难以回溯等问题，梳理“评价问题—逻辑关系—检查条件”的数据结构，推进\textbf{审核状态、修改建议、结果导出及 PDF 证据回溯}原型；结合实际案例开展字段校验与交互验收，形成可复用的专家复核流程。
  \item \textbf{垂域评测体系建设}：面向美股公司研究中的信息理解、财务推导与复杂推理任务，参与构建 \textbf{Rubrics}，拆解专家认知，明确证据要求、评分权重与重复扣分边界。
  \item \textbf{模型效果分析与优化}：结合 \textbf{LLM-as-a-Judge}、程序化财务检查和多轮重复实验开展 \textbf{Failure Case} 分析，区分内容遗漏、证据提取失败、口径冲突与评审误判；定位指标归一化、跨期匹配和高权重扣分误触发等问题，为数据规则、Prompt 与评测流程迭代提供依据。
  \item \textbf{数据治理}：参与建设 \textbf{1,700} 余份研报的资料储备，完成元数据整理、分类映射、附件核对与异常记录。
\end{itemize}

\ResumeItem{\textbf{产品经理实习生}}
[Tara]
[2026.07]
\begin{itemize}
  \item \textbf{产品设计}：负责 \textbf{AI 社交产品}用户端功能设计，撰写 \textbf{PRD}，完成用户卡片体系与信息架构设计。
  \item \textbf{原型开发}：使用 \textbf{AI 编程工具}独立完成\textbf{微信小程序}前端交互开发，推动产品从概念落地为可交互原型。
  \item \textbf{增长与 BD}：策划线下活动，探索垂直社群获客；跟踪 AI 社交赛道竞品，对接潜在合作方推进早期商务合作。
\end{itemize}

\ResumeItem{\textbf{投资者关系实习生}}
[奇绩创坛]
[2026.04—2026.06]
\begin{itemize}
  \item \textbf{投研支持}：研究 \textbf{Agentic AI} 领域发展态势，为陆奇博士的研究提供支持；对市场进行 mapping，梳理投资机会。
  \item \textbf{创业者网络}：链接 \textbf{500+} 创业者，对其学术禀赋与产业背景进行综合评估，定期跟踪动态并协助市场研究与数据分析，系统性研究潜在投资标的，及时发现新兴投资机会。
  \item \textbf{生态拓展}：链接 \textbf{50+} 海内外高校研究人员，针对视频生成、Agentic AI 等前沿领域通过深度访谈、行业调研、桌面研究快速开展研究，挖掘算力需求方并提供算力支持。
\end{itemize}

\ResumeItem{\textbf{投资银行部实习生}}
[中信证券]
[2026.01—2026.04]
\begin{itemize}
  \item \textbf{行业洞察}：参与某具身智能企业港股 \textbf{IPO} 项目，撰写招股书“业务”章节并完成同业竞对研究。
  \item \textbf{AI 提效}：使用 \textbf{Claude Code} 完成底稿核对、格式排版、资料整理等重复性任务，显著提升交付效率。
\end{itemize}

\section{项目}

\ResumeItem{\textbf{鉴真盾}}
[Qwen3.6、Agent、Prompt Engineering、Fine-tuning]
[2025.11]
\begin{itemize}
  \item \textbf{项目概况}：开发“鉴真盾”——基于\textbf{多模态大模型}的金融反欺诈人像风险识别系统，面向小微企业贷款预审场景，解决黑产中介伪造身份、批量攻击等传统信审痛点，通过人像识别判定申请人真实性风险并输出量化评分。
  \item \textbf{技术实现}：基于 \textbf{Qwen3.6} 模型完成数据标注、\textbf{Agent} 协作架构设计、模型测试与 \textbf{Prompt Engineering} 调优，并对模型进行 \textbf{Fine-tune}；建立中介欺诈特征库，精准识别六大类欺诈特征，负责全栈开发。
  \item \textbf{量化成果}：模型在测试集准确率达 \textbf{92\%}，实现风险量化评分机制；项目获“三创赛”\textbf{北京市一等奖}、“工行杯”全国大学生金融科技创新大赛\textbf{全国三等奖}。
\end{itemize}

\section{校园组织经历}

\ResumeItem{\textbf{UIBE AI 创协}}
[Co-founder]
[2025.11—至今]
\begin{itemize}
  \item \textbf{组织搭建}：从 0 到 1 发起 AI 创协，搭建校园运营、Startup 合作与 AI 应用三大板块，沉淀 \textbf{200+} 活跃用户。
  \item \textbf{SOP 构建}：面向 AI 企业设计校园增长合作模式 SOP，通过“校园运营—定制服务—数据复盘”的闭环，为企业提供校园增长解决方案；合作企业与生态覆盖智谱、Second Me、PandaAI、HyperKnow、观猹、知乎等。
  \item \textbf{营销实验}：调研学生群体用户需求并制定运营方案，通过线下活动、社媒、社群等多渠道投放与互动。
\end{itemize}

\section{能力}
\begin{itemize}
  \item \textbf{AI 工具与开发范式}：熟练使用 \textbf{Codex} 与 \textbf{Claude Code}（三个月 token 耗量近 \textbf{60 亿}）进行 \textbf{Vibe Coding} 开发；具备较强学习能力与 \textbf{AI Native} 的思考与协作方式，能快速上手并驾驭新兴 AI 开发工具链。
  \item \textbf{编程与数据分析}：熟练掌握 \textbf{Python}，擅长数据分析；了解 \textbf{SQL}、\textbf{C++}、\textbf{R}，对现代软件开发流程有基本认知。
\end{itemize}

\end{document}
